\section{Feasible-set universality and exact algebraic voltages}
\label{sec:algebraic}

The reduction preserves more than a yes-or-no answer. Its rational
solution-set correspondence transfers both topology and fields of definition.
We call two real sets \emph{rationally equivalent} if there is a
homeomorphism between them whose coordinate functions, and those of its
inverse, are rational functions with rational coefficients, with denominators
nonzero on their respective domains.
A set is \emph{basic closed over $\Q$} if it is described by a finite
conjunction of polynomial equalities and weak inequalities with rational
coefficients. General semialgebraic descriptions may also use disjunctions
and strict inequalities. This distinction matters for rational equivalence.

\begin{lemma}[Bounded arithmetic for basic closed sets]
\label{lem:basic-arithmetic}
Every compact basic closed set $S\subseteq\R^d$ over $\Q$ is rationally
equivalent to the solution set of a bounded $\ETRINV$ instance.
The equivalence can retain, for every original coordinate $t_i$, a designated
replacement coordinate $x_{j(i)}=1+\varepsilon t_i$, where
$\varepsilon>0$ is rational. All added coordinates are uniquely determined
rational functions of the original coordinates.
\end{lemma}

Appendix~\ref{app:arithmetic} proves the lemma using the conjunction-only
arithmetic constructions underlying Lemmas C--G of
\citet{AbrahamsenMiltzow2019}. It gives the gate identities, checks their
ranges, and avoids Boolean preprocessing. We assert existence here, without
a size bound for an arbitrary description of $S$.

\begin{theorem}
\label{thm:universality}
A compact semialgebraic set defined over $\Q$ is rationally equivalent
to the feasible voltage set of a rational $\RPF$ instance if and only if
it is basic closed over $\Q$.
For every prescribed fixed girth bound, the instance can satisfy all the
structural and finite-data restrictions of Theorem~\ref{thm:structural}.
\end{theorem}
\begin{proof}
For sufficiency, compose Lemma~\ref{lem:basic-arithmetic} with
Proposition~\ref{prop:rational-bijection}, or with the stronger construction
of Theorem~\ref{thm:structural}. Composition preserves rationality and
continuity in both directions. Necessity follows from
Proposition~\ref{prop:basic-invariance} below, because every rational RPF
voltage set is basic closed over $\Q$. Empty sets are included.
\end{proof}

\begin{proposition}[Basic closedness is preserved]
\label{prop:basic-invariance}
If a compact set $S$ is rationally equivalent over $\Q$ to a basic closed
set $T$ over $\Q$, then $S$ is basic closed over $\Q$.
\end{proposition}
\begin{proof}
The empty case is immediate. Otherwise write the equivalence and its
inverse as $F:S\to T$ and $G:T\to S$. Collect the polynomial denominators
of the coordinates of $F$. For each denominator of $G$, substitute $F$
and clear the forward denominators, without cancellation, to obtain a
polynomial numerator. Add these numerators to the collection. Every
collected polynomial is nonzero on $S$. Compactness supplies a rational
$\eta>0$ smaller than all their absolute values on $S$.

Let $S'$ consist of all ambient points $t$ satisfying three conditions:
every collected polynomial has square at least $\eta^2$;
$F(t)\in T$; and $G(F(t))=t$. The first condition makes all expressions
in the other two well defined. Clearing their denominators by positive
even powers turns them into polynomial equalities and weak inequalities
over $\Q$. Hence $S'$ is basic closed. The inverse identities give
$S\subseteq S'$. Conversely, if $t\in S'$, then $F(t)\in T$ and
$t=G(F(t))\in S$. Thus $S'=S$.
\end{proof}

The restriction in Theorem~\ref{thm:universality} cannot be dropped.
Consider the compact three-quadrant set
\begin{equation}
\label{eq:three-quadrants}
 S=[-1,1]^2\cap\bigl(\{(x,y):x\ge0\}\cup\{(x,y):y\ge0\}\bigr).
\end{equation}
It is not even locally basic closed at the origin. To see this, suppose
finitely many polynomial weak inequalities describe it near the origin;
equalities can be replaced by pairs of inequalities. Every nonzero
defining polynomial $p$ with $p(0)=0$ has a lowest nonzero homogeneous
part $H$ that is nonnegative in the three included open quadrants.
Its degree cannot be odd: the second and fourth quadrants are negatives
of each other, which would force $H$ to vanish on an open set.
Its degree is therefore even, so $H$ is nonnegative in the
remaining quadrant as well. Choose a direction in the excluded lower-left
quadrant avoiding the zero sets of these finitely many nonzero homogeneous
polynomials. All their values there are strictly positive. Along a
sufficiently short segment of that ray, every $p$ is positive; polynomials
with positive constant term also remain positive, and zero polynomials
impose no condition. This contradicts exclusion of that segment.

\begin{theorem}[Topological universality]
\label{thm:topological-universality}
Every compact semialgebraic set is semialgebraically homeomorphic to the
voltage feasible set of a rational RPF instance satisfying all restrictions
of Theorem~\ref{thm:structural}, for any prescribed fixed girth bound.
The homeomorphism need not be rational.
\end{theorem}
\begin{proof}
The empty case follows from Theorem~\ref{thm:universality}. A nonempty
compact semialgebraic set has a semialgebraic triangulation by a finite
simplicial complex $K$; see \citet[Theorem 1.1 and Section 1.2]{OhmotoShiota2017}.
Let its vertices be $1,\ldots,N$. The standard simplex realization of $K$ is
\begin{equation}
\label{eq:simplex-nonfaces}
 T_K=\left\{z\in\R^N:
 z_i\ge0,\quad \sum_i z_i=1,\quad
 \prod_{i\in F}z_i=0\ \text{for every nonface }F\subseteq\{1,\ldots,N\}
 \right\}.
\end{equation}
The support of $z$ is a face exactly when all the nonface equations hold.
Thus $T_K$ is the union of the standard simplices indexed by $K$.
The barycentric map from these standard simplices to any geometric
realization of $K$ is a piecewise linear homeomorphism: it is affine and
bijective on each simplex, and these maps agree on common faces.
Consequently $T_K$ is semialgebraically homeomorphic to the original set.
It is compact and basic closed over $\Q$, so
Theorem~\ref{thm:universality} applies. No bound on the size of the
triangulation or the list of nonfaces is asserted.
\end{proof}

The structured voltage sets therefore need not be connected, contractible,
or zero-dimensional. The equal-angle AC transfers preserve these magnitude
sets. Fixing a reference angle in each connected component gives the same
conclusions for rectangular voltages in the bus-angle-box model. Rational
equivalence retains algebraic information that a semialgebraic
homeomorphism can change; this is why the two universality statements
have different scopes.

\begin{theorem}
\label{thm:algebraic-degree}
Let $\alpha$ be a real algebraic number. There is a rational $\RPF$
instance with a unique feasible voltage vector $v$ such that every $v_i$
belongs to $\Q(\alpha)$ and $\Q(v_j)=\Q(\alpha)$ for some designated bus
$j$. Consequently, a uniquely feasible rational network can have a designated
bus voltage of any prescribed algebraic degree over $\Q$.
All structural restrictions of Theorem~\ref{thm:structural} can be imposed.
The same statement holds for unique AC magnitudes, and for unique rectangular
voltages in the reference-fixed bus-angle-box model.
\end{theorem}
\begin{proof}
First suppose $\alpha\notin\Q$. Choose a nonzero polynomial in $\Q[t]$
having $\alpha$ as a root and rational endpoints $a<b$ whose interval
contains $\alpha$ and no other root. Its root condition together with
$a\le t\le b$ defines the singleton $\{\alpha\}$.
Lemma~\ref{lem:basic-arithmetic} gives a bounded $\ETRINV$
solution set rationally equivalent to this singleton, so it too is a
singleton, say $\{x\}$. Rationality of the forward map gives
$x_i\in\Q(\alpha)$ for every coordinate.

The designated-coordinate conclusion of the lemma recovers an original
coordinate by a rational affine rescaling of one replacement coordinate. Thus
$\alpha=A x_j+B$ for some $A,B\in\Q$. Since $\alpha$ is irrational,
$A\ne0$, and hence $\Q(x_j)=\Q(\alpha)$.
The network construction has a unique extension of $x$, with rational
coordinate functions, and it retains $x_j$ as the voltage at its original
variable root. The planarization, connectors, and subdivisions retain that
root coordinate and add only uniquely determined rational functions of the
old coordinates. This proves the claims for irrational $\alpha$.

For rational $\alpha$, a single voltage-$1$, zero-injection bus is sufficient:
its unique voltage generates $\Q=\Q(\alpha)$ and its graph meets every
stated restriction. To obtain degree $d\ge2$, take the positive real root
of $t^d-2$; this polynomial is irreducible by Eisenstein's criterion at $2$.
Degree one is the preceding rational example. Finally, the exact AC transfers
preserve precisely the RPF magnitude set, and reference-fixed angle boxes
force every phase to zero.
\end{proof}

The designated-coordinate assertion uses the affine recovery property,
not merely preservation of the field generated by all coordinates. The
theorem is an existence statement with no size bound on the affine recovery
coefficients. It does not by itself exclude polynomially verifiable
symbolic certificates; the obstruction to such certificates is the
complexity consequence in Corollary~\ref{cor:rpf-np}.
